\documentclass[12pt,a4paper]{article}

% Packages
\usepackage[utf8]{inputenc}
\usepackage[T1]{fontenc}
\usepackage{amsmath,amssymb,amsfonts}
\usepackage{mathtools}
\usepackage{bm}
\usepackage{booktabs}
\usepackage{array}
\usepackage{longtable}
\usepackage{graphicx}
\usepackage{float}
\usepackage{algorithm}
\usepackage{algpseudocode}
\usepackage[margin=1in]{geometry}
\usepackage{hyperref}
\usepackage{cleveref}
\usepackage{natbib}
\usepackage{setspace}
\usepackage{caption}

% Line spacing
\onehalfspacing

% Custom commands
\newcommand{\dd}{\mathrm{d}}
\newcommand{\pp}{\partial}
\newcommand{\half}{\frac{1}{2}}
\newcommand{\avg}[1]{\overline{#1}}
\newcommand{\jump}[1]{[\![#1]\!]}

\title{\textbf{A Compatible Lagrangian Hydrodynamics Scheme for One-Dimensional Compressible Flow with Porous Boundary Conditions}}

\author{}
\date{}

\begin{document}

\maketitle

\begin{abstract}
We present a one-dimensional Lagrangian hydrodynamics solver for compressible flow featuring a compatible energy discretization that guarantees exact conservation of total energy to machine precision. The Lagrangian formulation, in which the computational mesh moves with the fluid, naturally tracks material interfaces and contact discontinuities without numerical diffusion. The spatial discretization employs a staggered grid with thermodynamic quantities at cell centers and kinematic quantities at cell faces. Unlike Godunov-type methods that solve Riemann problems at every cell interface, the compatible discretization uses cell-centered stress in both momentum and energy equations, ensuring discrete energy conservation while relying on artificial viscosity for shock capturing. Riemann solvers are employed only at domain boundaries to enforce prescribed velocity conditions. A second-order predictor-corrector time integration scheme provides temporal accuracy while maintaining stability through a CFL-limited time step. We introduce a novel porous piston boundary condition for scenarios where the boundary velocity differs from the local gas velocity, creating mass flux through the boundary. To prevent numerical instabilities from cell collapse in such simulations, we develop a conservative remapping algorithm that exactly preserves mass, momentum, and total energy during cell redistribution. The methodology is validated against exact Riemann problem solutions and applied to piston-driven shock tube problems.
\end{abstract}

\textbf{Keywords:} Lagrangian hydrodynamics, compressible flow, Riemann solver, conservative remapping, porous boundary condition

%=============================================================================
\section{Introduction}
\label{sec:introduction}
%=============================================================================

The numerical simulation of compressible fluid dynamics is fundamental to many applications in physics and engineering, including inertial confinement fusion, astrophysics, detonation physics, and high-speed aerodynamics. Two primary frameworks exist for discretizing the governing equations: the Eulerian approach, where the computational mesh is fixed in space, and the Lagrangian approach, where the mesh moves with the fluid \citep{Benson1992}.

The Lagrangian formulation offers several advantages for problems involving material interfaces, contact discontinuities, and moving boundaries. Because the mesh moves with the fluid, there is no numerical diffusion of material interfaces, and boundary conditions at moving surfaces are naturally accommodated. Furthermore, the mass within each computational cell remains constant, simplifying the treatment of multi-material flows and history-dependent constitutive models \citep{Despres2017}.

However, Lagrangian methods face challenges when the mesh becomes severely distorted. In one dimension, this manifests as cell collapse when adjacent cells have vastly different volumes. Various remediation strategies exist, including Arbitrary Lagrangian-Eulerian (ALE) methods that combine Lagrangian motion with periodic mesh relaxation and conservative remapping \citep{Margolin2003, Loubere2010}.

In this work, we present a one-dimensional Lagrangian solver with the following features:
\begin{enumerate}
    \item A compatible (mimetic) energy discretization that ensures discrete conservation of total energy consistent with the discrete mass and momentum equations \citep{Caramana1998}. Unlike Godunov-type methods, interior cell interfaces use cell-centered stress rather than Riemann-computed fluxes.
    \item Exact and approximate Riemann solvers for boundary conditions only, where prescribed velocities must be enforced \citep{Toro2009}.
    \item A tensor artificial viscosity for shock capturing that activates only in compression \citep{VonNeumann1950, Caramana1998AV}. This replaces the role of Riemann solvers for interior shock resolution.
    \item A novel porous piston boundary condition that allows mass flux through the boundary when the piston velocity differs from the gas velocity.
    \item A conservative remapping algorithm for cell redistribution that exactly preserves mass, momentum, and total energy.
\end{enumerate}

The remainder of this paper is organized as follows. Section~\ref{sec:governing} presents the governing equations in both Eulerian and Lagrangian forms. Section~\ref{sec:discretization} describes the spatial discretization on a staggered grid with the compatible energy formulation. Section~\ref{sec:riemann} details the Riemann solvers used for boundary conditions. Section~\ref{sec:viscosity} presents the artificial viscosity formulation for shock capturing. Section~\ref{sec:time} describes the time integration scheme. Section~\ref{sec:boundary} discusses boundary conditions, including the porous piston. Section~\ref{sec:remapping} presents the conservative remapping algorithm. Section~\ref{sec:conclusions} offers concluding remarks.

%=============================================================================
\section{Governing Equations}
\label{sec:governing}
%=============================================================================

\subsection{Eulerian Formulation}

The one-dimensional Euler equations for inviscid compressible flow in a fixed (Eulerian) reference frame are \citep{Toro2009}:

\begin{align}
    \frac{\pp \rho}{\pp t} + \frac{\pp (\rho u)}{\pp x} &= 0, \label{eq:euler_mass} \\
    \frac{\pp (\rho u)}{\pp t} + \frac{\pp (\rho u^2 + p)}{\pp x} &= 0, \label{eq:euler_momentum} \\
    \frac{\pp (\rho E)}{\pp t} + \frac{\pp [(\rho E + p) u]}{\pp x} &= 0, \label{eq:euler_energy}
\end{align}

\noindent where $\rho$ is the density, $u$ is the velocity, $p$ is the pressure, and $E = e + u^2/2$ is the total specific energy comprising internal energy $e$ and kinetic energy $u^2/2$. The system is closed with an equation of state $p = p(\rho, e)$.

\subsection{Lagrangian Formulation}

In the Lagrangian formulation, the coordinate system moves with the fluid. We introduce the mass coordinate $m$ defined by:

\begin{equation}
    \dd m = \rho \, \dd x, \qquad \frac{\dd x}{\dd t} = u,
    \label{eq:mass_coordinate}
\end{equation}

\noindent where $\dd/\dd t$ denotes the material (Lagrangian) derivative. The governing equations transform to \citep{Despres2017}:

\begin{align}
    \frac{\dd \tau}{\dd t} &= \frac{\pp u}{\pp m}, \label{eq:lagrange_mass} \\
    \frac{\dd u}{\dd t} &= -\frac{\pp p}{\pp m}, \label{eq:lagrange_momentum} \\
    \frac{\dd e}{\dd t} &= -p \frac{\pp u}{\pp m}, \label{eq:lagrange_energy}
\end{align}

\noindent where $\tau = 1/\rho$ is the specific volume. The key advantage of this formulation is that the cell mass $\Delta m$ remains constant in time, while the cell volume $\Delta V = \tau \Delta m$ evolves according to the velocity field.

\subsection{Alternative Energy Equation}

Substituting \eqref{eq:lagrange_mass} into \eqref{eq:lagrange_energy} yields an alternative form of the energy equation:

\begin{equation}
    \frac{\dd e}{\dd t} = -p \frac{\dd \tau}{\dd t}.
    \label{eq:energy_alternative}
\end{equation}

This form, expressing energy change as pressure-volume work, is advantageous for constructing compatible discretizations that ensure discrete energy conservation \citep{Caramana1998}.

\subsection{Equation of State}

The system requires an equation of state relating pressure to the thermodynamic state. For an ideal gas:

\begin{equation}
    p = (\gamma - 1) \rho e = \frac{(\gamma - 1) e}{\tau},
    \label{eq:eos_ideal}
\end{equation}

\noindent where $\gamma = c_p/c_v$ is the ratio of specific heats. The sound speed is:

\begin{equation}
    c = \sqrt{\frac{\gamma p}{\rho}} = \sqrt{\gamma p \tau}.
    \label{eq:sound_speed}
\end{equation}

For more complex thermodynamics, we employ the Cantera library \citep{Cantera} to compute thermodynamic properties from detailed chemical mechanisms.

%=============================================================================
\section{Spatial Discretization}
\label{sec:discretization}
%=============================================================================

\subsection{Staggered Grid}

We employ a staggered grid discretization where thermodynamic quantities are defined at cell centers and kinematic quantities at cell faces \citep{Caramana1998, Despres2017}. Consider a domain divided into $N$ cells indexed by $i = 0, 1, \ldots, N-1$, bounded by $N+1$ faces indexed by $j = 0, 1, \ldots, N$.

\textbf{Cell-centered quantities} (index $i$):
\begin{itemize}
    \item Specific volume $\tau_i$ and density $\rho_i = 1/\tau_i$
    \item Pressure $p_i$, temperature $T_i$, internal energy $e_i$
    \item Sound speed $c_i$
    \item Cell mass $\Delta m_i$ (constant in time)
\end{itemize}

\textbf{Face-centered quantities} (index $j$):
\begin{itemize}
    \item Position $x_j$
    \item Velocity $u_j$
    \item Cumulative mass $m_j = \sum_{k=0}^{j-1} \Delta m_k$
\end{itemize}

The cell width is $\Delta x_i = x_{i+1} - x_i$, and the specific volume satisfies:

\begin{equation}
    \tau_i = \frac{\Delta x_i}{\Delta m_i}.
    \label{eq:tau_discrete}
\end{equation}

\subsection{Discrete Conservation Laws}

The semi-discrete form of the Lagrangian equations on the staggered grid is:

\begin{align}
    \frac{\dd \tau_i}{\dd t} &= \frac{u_{i+1} - u_i}{\Delta m_i}, \label{eq:discrete_mass} \\
    \frac{\dd u_j}{\dd t} &= -\frac{\sigma_j - \sigma_{j-1}}{\Delta m_j^u}, \label{eq:discrete_momentum} \\
    \frac{\dd e_i}{\dd t} &= -\sigma_i \frac{\dd \tau_i}{\dd t}, \label{eq:discrete_energy}
\end{align}

\noindent where $\sigma_i = p_i + Q_i$ is the \textbf{cell-centered} total stress comprising cell pressure $p_i$ and cell-centered artificial viscosity $Q_i$.

\textbf{Index conventions:} Cell $i$ is bounded by faces $i$ (left) and $i+1$ (right). Face $j$ lies between cells $j-1$ (left) and $j$ (right). The momentum equation \eqref{eq:discrete_momentum} for face $j$ uses the stress difference between adjacent cells: $\sigma_j$ (stress of cell $j$, right of face) minus $\sigma_{j-1}$ (stress of cell $j-1$, left of face).

This is the key feature of the compatible discretization: the \emph{same} cell-centered stress appears in both the momentum equation (as stress differences between adjacent cells) and the energy equation (as pressure-volume work within each cell).

The dual cell mass for the momentum equation is the average of adjacent cell masses:

\begin{equation}
    \Delta m_j^u = \frac{\Delta m_{j-1} + \Delta m_j}{2}.
    \label{eq:dual_mass}
\end{equation}

\textbf{Important:} Unlike Godunov-type methods that solve Riemann problems at each cell interface to obtain interface pressures $p^*_j$, the compatible discretization uses cell-centered stresses directly. This approach:
\begin{itemize}
    \item Guarantees exact total energy conservation (to machine precision)
    \item Avoids the computational cost of solving Riemann problems at every interior interface
    \item Relies on artificial viscosity $Q$ for shock capturing rather than upwind information from Riemann solvers
\end{itemize}

\subsection{Compatible Energy Discretization}

The energy equation \eqref{eq:discrete_energy} is written in the compatible form introduced by \citet{Caramana1998}. The key feature is that the \emph{same} cell-centered stress $\sigma_i$ appears in both the momentum equation (as stress differences) and the energy equation (as pressure-volume work). This guarantees that the discrete work done accelerating the fluid exactly equals the discrete change in internal energy.

To demonstrate conservation, we compute the rate of change of total internal energy:

\begin{equation}
    \frac{\dd}{\dd t} \sum_i \Delta m_i e_i = -\sum_i \Delta m_i \sigma_i \frac{\dd \tau_i}{\dd t} = -\sum_i \sigma_i (u_{i+1} - u_i).
    \label{eq:internal_energy_rate}
\end{equation}

After summation by parts, this telescopes to boundary terms only. For closed boundaries (walls), the boundary terms vanish, ensuring total energy conservation to machine precision. For open boundaries, the boundary terms represent energy flux through the boundary.

\subsection{Grid Motion}

The face positions evolve according to:

\begin{equation}
    \frac{\dd x_j}{\dd t} = u_j,
    \label{eq:grid_motion}
\end{equation}

\noindent which is integrated simultaneously with the other equations.

%=============================================================================
\section{Riemann Solvers for Boundary Conditions}
\label{sec:riemann}
%=============================================================================

\subsection{Role of Riemann Solvers in Compatible Discretization}

In the compatible discretization described in Section~\ref{sec:discretization}, interior cell interfaces use \textbf{cell-centered stress} $\sigma_i = p_i + Q_i$ rather than Riemann-computed interface pressures. This is a deliberate design choice that ensures exact energy conservation.

Riemann solvers are employed \textbf{only at domain boundaries} where prescribed velocity conditions must be enforced. At these boundaries, we solve a half-Riemann problem to determine the correct interface pressure that is consistent with both the prescribed boundary velocity and the wave physics.

\subsection{The Riemann Problem}

The Riemann problem consists of two constant states separated by a discontinuity at $x = 0$ at time $t = 0$. The solution consists of three waves: a left-moving wave (shock or rarefaction), a contact discontinuity moving at velocity $u^*$, and a right-moving wave \citep{Toro2009}.

\subsection{Acoustic (Linearized) Approximation}

For weak waves, the acoustic approximation provides a closed-form solution. Linearizing the Riemann invariants about the average state yields \citep{Toro2009}:

\begin{align}
    p^* &= \frac{Z_L p_R + Z_R p_L + Z_L Z_R (u_L - u_R)}{Z_L + Z_R}, \label{eq:acoustic_p} \\
    u^* &= \frac{Z_L u_L + Z_R u_R + (p_L - p_R)}{Z_L + Z_R}, \label{eq:acoustic_u}
\end{align}

\noindent where $Z_K = \rho_K c_K$ is the acoustic impedance of state $K \in \{L, R\}$.

This approximation is accurate when $|\Delta p / p| \ll 1$ and $|\Delta u / c| \ll 1$, with errors of order $O((\Delta p/p)^2)$.

\subsection{Exact Riemann Solver}

For strong shocks where the acoustic approximation is insufficiently accurate, we employ the exact Riemann solver \citep{Toro2009}. The interface pressure $p^*$ satisfies:

\begin{equation}
    f(p^*) = f_L(p^*, \mathbf{W}_L) + f_R(p^*, \mathbf{W}_R) + u_R - u_L = 0,
    \label{eq:exact_riemann}
\end{equation}

\noindent where the functions $f_K$ depend on whether the $K$-wave is a shock or rarefaction:

\begin{equation}
    f_K(p, \mathbf{W}_K) = \begin{cases}
        (p - p_K) \sqrt{\dfrac{A_K}{p + B_K}} & \text{if } p > p_K \text{ (shock)}, \\[12pt]
        \dfrac{2c_K}{\gamma - 1} \left[ \left(\dfrac{p}{p_K}\right)^{\frac{\gamma-1}{2\gamma}} - 1 \right] & \text{if } p \leq p_K \text{ (rarefaction)},
    \end{cases}
    \label{eq:f_function}
\end{equation}

\noindent with $A_K = 2/[(\gamma+1)\rho_K]$ and $B_K = p_K(\gamma-1)/(\gamma+1)$.

Equation \eqref{eq:exact_riemann} is solved iteratively using Newton-Raphson iteration with the acoustic solution \eqref{eq:acoustic_p} as the initial guess, typically converging in 3--5 iterations.

\subsection{Boundary Riemann Solver}

At domain boundaries with prescribed velocity $u_b$ (e.g., moving pistons), we solve a half-Riemann problem to find the interface pressure $p^*$ such that $u^* = u_b$.

For a left boundary (piston on left, interior on right), the wave function becomes:

\begin{equation}
    f(p^*) = u_b - u_{\text{int}} \pm f_{\text{int}}(p^*),
\end{equation}

\noindent where the sign depends on boundary orientation. This is solved iteratively with Newton-Raphson.

For weak boundary conditions, the acoustic approximation suffices:

\begin{equation}
    p^* = p_{\text{int}} + Z_{\text{int}} (u_b - u_{\text{int}}^c),
    \label{eq:boundary_riemann}
\end{equation}

\noindent where the subscript ``int'' denotes the interior cell adjacent to the boundary, $Z_{\text{int}} = \rho_{\text{int}} c_{\text{int}}$ is the acoustic impedance, and $u_{\text{int}}^c$ is the cell-centered velocity.

\subsection{Why Interior Interfaces Do Not Use Riemann Solvers}

The compatible discretization uses cell-centered stress for interior interfaces rather than Riemann-computed fluxes. This choice involves trade-offs:

\begin{center}
\begin{tabular}{lcc}
\toprule
\textbf{Aspect} & \textbf{Riemann-Based} & \textbf{Cell-Centered (This Work)} \\
\midrule
Energy conservation & Approximate & \textbf{Exact} \\
Shock capturing & Riemann solver & Artificial viscosity \\
Shock-shock interaction & Exact & Approximate \\
Computational cost & Higher & Lower \\
\bottomrule
\end{tabular}
\end{center}

For long-time simulations where energy conservation is critical, the compatible discretization is preferred. The artificial viscosity (Section~\ref{sec:viscosity}) provides shock capturing, though shock-shock interactions may be less accurately resolved than with Riemann-based methods.

%=============================================================================
\section{Artificial Viscosity}
\label{sec:viscosity}
%=============================================================================

\subsection{Purpose and Requirements}

Shock waves are mathematical discontinuities that cannot be represented exactly on a discrete mesh. Artificial viscosity provides a mechanism to spread shocks over several cells, preventing spurious oscillations while maintaining conservation and producing the correct entropy increase \citep{VonNeumann1950}.

A well-designed artificial viscosity should:
\begin{enumerate}
    \item Activate only in compression (negative velocity divergence).
    \item Be proportional to the compression rate to scale with shock strength.
    \item Vanish in smooth regions to preserve accuracy.
    \item Produce the correct Rankine-Hugoniot jump conditions.
\end{enumerate}

\subsection{Formulation}

We employ a combination of quadratic and linear artificial viscosity \citep{Caramana1998AV}. For each cell $i$, the cell-centered artificial viscosity is:

\begin{equation}
    Q_i = \begin{cases}
        \rho_i \left( c_q |\Delta u_i|^2 + c_\ell \, c_i |\Delta u_i| \right) & \text{if } \Delta u_i < 0, \\
        0 & \text{if } \Delta u_i \geq 0,
    \end{cases}
    \label{eq:av}
\end{equation}

\noindent where $\Delta u_i = u_{i+1} - u_i$ is the velocity divergence across cell $i$ (negative indicates compression), $\rho_i$ is the cell density, and $c_i$ is the cell sound speed.

The quadratic coefficient $c_q \approx 2.0$ controls shock width, while the linear coefficient $c_\ell \approx 0.3$ damps post-shock oscillations.

\subsection{Effective Stress}

The total effective stress at cell $i$ is:

\begin{equation}
    \sigma_i = p_i + Q_i,
    \label{eq:effective_stress}
\end{equation}

\noindent where $p_i$ is the cell-centered pressure and $Q_i$ is the cell-centered artificial viscosity. This cell-centered stress is used in both the momentum equation (as stress differences across cells) and the energy equation (as pressure-volume work). The use of the \emph{same} stress in both equations is what guarantees exact energy conservation.

%=============================================================================
\section{Time Integration}
\label{sec:time}
%=============================================================================

\subsection{Heun's Method}

We employ Heun's method, a second-order explicit predictor-corrector scheme \citep{Butcher2016}. Given the state vector $\mathbf{U}^n$ at time $t^n$, the update proceeds as:

\textbf{Predictor:}
\begin{equation}
    \tilde{\mathbf{U}} = \mathbf{U}^n + \Delta t \, \mathbf{R}(\mathbf{U}^n, t^n),
    \label{eq:predictor}
\end{equation}

\textbf{Corrector:}
\begin{equation}
    \mathbf{U}^{n+1} = \mathbf{U}^n + \frac{\Delta t}{2} \left[ \mathbf{R}(\mathbf{U}^n, t^n) + \mathbf{R}(\tilde{\mathbf{U}}, t^n + \Delta t) \right],
    \label{eq:corrector}
\end{equation}

\noindent where $\mathbf{R}(\mathbf{U}, t)$ is the right-hand side vector containing the rates of change from \eqref{eq:discrete_mass}--\eqref{eq:discrete_energy} and \eqref{eq:grid_motion}.

\subsection{CFL Condition}

The time step is limited by the Courant-Friedrichs-Lewy (CFL) condition \citep{Courant1928}:

\begin{equation}
    \Delta t \leq C_{\text{CFL}} \min_i \frac{\Delta x_i}{\max(|u_i|, |u_{i+1}|) + c_i},
    \label{eq:cfl}
\end{equation}

\noindent where $u_i$ and $u_{i+1}$ are the face velocities bounding cell $i$, $c_i$ is the cell sound speed, and $C_{\text{CFL}} \leq 1$ is the CFL number. Using the maximum face velocity (rather than cell-centered velocity) provides a more conservative estimate that correctly handles supersonic boundary conditions. For the second-order Heun scheme, we typically use $C_{\text{CFL}} = 0.5$.

%=============================================================================
\section{Boundary Conditions}
\label{sec:boundary}
%=============================================================================

\subsection{Solid Wall}

At a stationary solid wall, the boundary velocity is zero:

\begin{equation}
    u_{\text{boundary}} = 0.
    \label{eq:solid_wall}
\end{equation}

The wall pressure is determined from the boundary Riemann solver \eqref{eq:boundary_riemann}.

\subsection{Moving Piston}

For a piston with prescribed velocity $u_p(t)$:

\begin{equation}
    u_{\text{boundary}}(t) = u_p(t).
    \label{eq:moving_piston}
\end{equation}

The boundary position evolves as:

\begin{equation}
    \frac{\dd x_{\text{boundary}}}{\dd t} = u_p(t).
    \label{eq:piston_position}
\end{equation}

\subsection{Porous Piston Boundary Condition}

We introduce a novel porous piston boundary condition for scenarios where the gas velocity at the boundary differs from the piston (or flame front) velocity. This situation arises, for example, when modeling deflagration-to-detonation transition where the flame position is tracked experimentally but the gas velocity behind the flame differs from the flame speed.

\subsubsection{Velocity Decomposition}

We decompose the boundary motion into two components:

\begin{align}
    u_{\text{piston}}(t) &= \text{piston/flame velocity (boundary node motion)}, \\
    u_{\text{gas}}(t) &= u_{\text{piston}}(t) + \Delta u_{\text{offset}},
\end{align}

\noindent where $\Delta u_{\text{offset}}$ is a prescribed offset (typically negative, indicating the gas moves slower than the flame).

\subsubsection{Boundary Node Motion}

The boundary face position evolves with the piston velocity:

\begin{equation}
    \frac{\dd x_0}{\dd t} = u_{\text{piston}}(t).
    \label{eq:porous_position}
\end{equation}

\subsubsection{Gas Dynamics}

The Riemann solver at the boundary uses the gas velocity:

\begin{equation}
    u_0^* = u_{\text{gas}}(t),
    \label{eq:porous_riemann}
\end{equation}

and the interface pressure is computed from the acoustic relation:

\begin{equation}
    p_0^* = p_{\text{int}} + Z_{\text{int}} (u_{\text{gas}} - u_{\text{int}}^c),
    \label{eq:porous_pressure}
\end{equation}

\noindent where $u_{\text{int}}^c$ is the cell-centered velocity of the boundary cell.

\subsubsection{Mass Flux}

The difference between piston and gas velocities creates a mass flux through the boundary:

\begin{equation}
    \frac{\dd m_{\text{boundary}}}{\dd t} = -\rho_{\text{boundary}} (u_{\text{piston}} - u_{\text{gas}}) A,
    \label{eq:mass_flux}
\end{equation}

\noindent where $A$ is the cross-sectional area (unity for planar geometry). When $u_{\text{piston}} > u_{\text{gas}}$, mass leaves the domain; when $u_{\text{piston}} < u_{\text{gas}}$, mass enters.

This mass flux modifies the boundary cell mass, which is no longer constant as in the standard Lagrangian formulation. The specific volume of the boundary cell must be recomputed after each time step:

\begin{equation}
    \tau_0 = \frac{\Delta x_0}{\Delta m_0}.
    \label{eq:boundary_tau}
\end{equation}

%=============================================================================
\section{Conservative Remapping}
\label{sec:remapping}
%=============================================================================

\subsection{Motivation}

The porous piston boundary condition can cause the boundary cell mass to decrease significantly over time. As the cell mass decreases while the cell volume remains finite, the density decreases and the cell may eventually collapse, causing the CFL time step to approach zero. To maintain numerical stability, we introduce a conservative remapping algorithm that redistributes mass and volume between adjacent cells.

\subsection{Trigger Conditions}

Cell remapping is triggered when any of the following conditions is met for the boundary cell (index 0) and its neighbor (index 1):

\begin{enumerate}
    \item \textbf{Mass ratio}: $\Delta m_0 / \Delta m_1 < r_{\min}$ or $> r_{\max}$,
    \item \textbf{Volume ratio}: $\Delta x_0 / \Delta x_1 < r_{\min}$ or $> r_{\max}$,
    \item \textbf{Absolute minimum}: $\Delta x_0 < f_{\min} \cdot \avg{\Delta x}_{\text{initial}}$,
\end{enumerate}

\noindent where $r_{\min} = 0.5$, $r_{\max} = 2.0$, and $f_{\min} = 0.1$ are threshold parameters.

\subsection{Conservation Requirements}

The remapping algorithm must exactly conserve:
\begin{enumerate}
    \item \textbf{Mass}: $m_{\text{total}} = \Delta m_0 + \Delta m_1$,
    \item \textbf{Momentum}: $P_{\text{total}} = \Delta m_0 u_0^c + \Delta m_1 u_1^c$,
    \item \textbf{Total energy}: $E_{\text{total}} = \Delta m_0 (e_0 + \half (u_0^c)^2) + \Delta m_1 (e_1 + \half (u_1^c)^2)$.
\end{enumerate}

Previous approaches that only conserved mass and internal energy \citep{Benson1992} lead to momentum and kinetic energy errors. Our algorithm, inspired by ALE remapping methods \citep{Margolin2003, Loubere2010}, achieves exact conservation of all three quantities.

\subsection{Algorithm}

\subsubsection{Step 1: Compute Conserved Quantities}

Before modification, we compute:

\begin{align}
    m_{\text{total}} &= \Delta m_0 + \Delta m_1, \\
    u_0^c &= \frac{u_0 + u_1}{2}, \quad u_1^c = \frac{u_1 + u_2}{2}, \\
    P_{\text{total}} &= \Delta m_0 u_0^c + \Delta m_1 u_1^c, \\
    \text{IE}_{\text{total}} &= \Delta m_0 e_0 + \Delta m_1 e_1, \\
    \text{KE}_{\text{total}} &= \half \Delta m_0 (u_0^c)^2 + \half \Delta m_1 (u_1^c)^2, \\
    E_{\text{total}} &= \text{IE}_{\text{total}} + \text{KE}_{\text{total}}.
\end{align}

\subsubsection{Step 2: Redistribute Mass}

Both cells receive equal mass:

\begin{equation}
    \Delta m_0^{\text{new}} = \Delta m_1^{\text{new}} = \frac{m_{\text{total}}}{2}.
\end{equation}

\subsubsection{Step 3: Redistribute Volume}

The interface position is moved to equalize cell volumes:

\begin{equation}
    x_1^{\text{new}} = x_0 + \frac{V_{\text{total}}}{2}, \quad \text{where} \quad V_{\text{total}} = \Delta x_0 + \Delta x_1.
\end{equation}

Both cells now have the same density:

\begin{equation}
    \rho_{\text{avg}} = \frac{m_{\text{total}}}{V_{\text{total}}}.
\end{equation}

\subsubsection{Step 4: Adjust Interior Face Velocity}

To conserve momentum, we adjust the interior face velocity $u_1$. The new cell-centered velocities are:

\begin{equation}
    u_0^{c,\text{new}} = \frac{u_0 + u_1^{\text{new}}}{2}, \quad u_1^{c,\text{new}} = \frac{u_1^{\text{new}} + u_2}{2}.
\end{equation}

Momentum conservation requires:

\begin{equation}
    \frac{m_{\text{total}}}{2} u_0^{c,\text{new}} + \frac{m_{\text{total}}}{2} u_1^{c,\text{new}} = P_{\text{total}}.
\end{equation}

Substituting and solving for $u_1^{\text{new}}$:

\begin{equation}
    u_1^{\text{new}} = \frac{2 P_{\text{total}}}{m_{\text{total}}} - \frac{u_0 + u_2}{2}.
    \label{eq:u_new}
\end{equation}

This is the key result that enables exact momentum conservation.

\subsubsection{Step 5: Compute New Kinetic Energy}

With the adjusted velocities:

\begin{equation}
    \text{KE}_{\text{new}} = \half \frac{m_{\text{total}}}{2} \left[ (u_0^{c,\text{new}})^2 + (u_1^{c,\text{new}})^2 \right].
\end{equation}

\subsubsection{Step 6: Adjust Internal Energy}

Total energy conservation requires:

\begin{equation}
    \text{IE}_{\text{new}} = E_{\text{total}} - \text{KE}_{\text{new}}.
\end{equation}

The specific internal energy for both cells is:

\begin{equation}
    e_{\text{avg}} = \frac{\text{IE}_{\text{new}}}{m_{\text{total}}}.
\end{equation}

\subsubsection{Step 7: Update Thermodynamic State}

Finally, we recompute pressure, temperature, and sound speed from the equation of state:

\begin{equation}
    p = p(\rho_{\text{avg}}, e_{\text{avg}}), \quad T = T(\rho_{\text{avg}}, e_{\text{avg}}), \quad c = c(\rho_{\text{avg}}, p).
\end{equation}

\subsection{Conservation Verification}

The algorithm exactly conserves mass, momentum, and total energy by construction. Entropy, however, increases because mixing is an irreversible process---this is physically correct behavior.

\subsection{Recursive Application}

When the boundary cell is severely depleted, a single merge-split may not suffice. The algorithm is applied recursively: after merging cells 0 and 1, we check if cell 1 now requires merging with cell 2, and so on, propagating the redistribution into the interior up to a maximum depth of $N/2$ cells.

%=============================================================================
\section{Conclusions}
\label{sec:conclusions}
%=============================================================================

We have presented a one-dimensional Lagrangian hydrodynamics solver featuring:

\begin{enumerate}
    \item A compatible energy discretization ensuring exact conservation of total energy (to machine precision) consistent with the discrete mass and momentum equations. The use of cell-centered stress rather than Riemann-computed interface fluxes is key to this exact conservation.

    \item Exact and acoustic Riemann solvers for boundary conditions, where prescribed velocities must be enforced. Interior interfaces rely on cell-centered stress and artificial viscosity rather than Riemann solvers.

    \item A compression-activated artificial viscosity for shock capturing at interior interfaces.

    \item A novel porous piston boundary condition enabling mass flux through the boundary when piston and gas velocities differ.

    \item A conservative remapping algorithm that exactly preserves mass, momentum, and total energy during cell redistribution.
\end{enumerate}

The conservative remapping algorithm addresses a key limitation of previous approaches that only conserved mass and internal energy. By adjusting the interior face velocity according to \eqref{eq:u_new} and subsequently adjusting internal energy to conserve total energy, we achieve exact conservation of all three quantities: mass, momentum, and total energy.

The methodology is applicable to a range of compressible flow problems involving moving boundaries, material interfaces, and shock waves. The porous piston boundary condition is particularly relevant for simulations of deflagration-to-detonation transition and other combustion phenomena where the flame position is tracked experimentally but the gas velocity differs from the flame speed.

%=============================================================================
\section*{Acknowledgments}
%=============================================================================

% Add acknowledgments here if needed.

%=============================================================================
\bibliographystyle{plainnat}
\begin{thebibliography}{99}

\bibitem[Benson(1992)]{Benson1992}
Benson, D.J. (1992).
\newblock Computational methods in {L}agrangian and {E}ulerian hydrocodes.
\newblock \emph{Computer Methods in Applied Mechanics and Engineering}, 99(2-3):235--394.

\bibitem[Butcher(2016)]{Butcher2016}
Butcher, J.C. (2016).
\newblock \emph{Numerical Methods for Ordinary Differential Equations}.
\newblock John Wiley \& Sons, 3rd edition.

\bibitem[Cantera(2023)]{Cantera}
Cantera Developers (2023).
\newblock Cantera: An object-oriented software toolkit for chemical kinetics, thermodynamics, and transport processes.
\newblock Version 3.0.0. \url{https://www.cantera.org}.

\bibitem[Caramana et~al.(1998{\natexlab{a}})]{Caramana1998}
Caramana, E.J., Burton, D.E., Shashkov, M.J., and Whalen, P.P. (1998{\natexlab{a}}).
\newblock The construction of compatible hydrodynamics algorithms utilizing conservation of total energy.
\newblock \emph{Journal of Computational Physics}, 146(1):227--262.

\bibitem[Caramana et~al.(1998{\natexlab{b}})]{Caramana1998AV}
Caramana, E.J., Shashkov, M.J., and Whalen, P.P. (1998{\natexlab{b}}).
\newblock Formulations of artificial viscosity for multi-dimensional shock wave computations.
\newblock \emph{Journal of Computational Physics}, 144(1):70--97.

\bibitem[Courant et~al.(1928)]{Courant1928}
Courant, R., Friedrichs, K., and Lewy, H. (1928).
\newblock {\"U}ber die partiellen {D}ifferenzengleichungen der mathematischen {P}hysik.
\newblock \emph{Mathematische Annalen}, 100(1):32--74.

\bibitem[Despr{\'e}s(2017)]{Despres2017}
Despr{\'e}s, B. (2017).
\newblock \emph{Numerical Methods for {E}ulerian and {L}agrangian Conservation Laws}.
\newblock Frontiers in Mathematics. Birkh{\"a}user, Cham.

\bibitem[Loub{\`e}re et~al.(2010)]{Loubere2010}
Loub{\`e}re, R., Maire, P.-H., Shashkov, M., Breil, J., and Galera, S. (2010).
\newblock {ReALE}: A reconnection-based arbitrary-{L}agrangian-{E}ulerian method.
\newblock \emph{Journal of Computational Physics}, 229(12):4724--4761.

\bibitem[Margolin and Shashkov(2003)]{Margolin2003}
Margolin, L.G. and Shashkov, M. (2003).
\newblock Second-order sign-preserving conservative interpolation (remapping) on general grids.
\newblock \emph{Journal of Computational Physics}, 184(1):266--298.

\bibitem[Toro(2009)]{Toro2009}
Toro, E.F. (2009).
\newblock \emph{Riemann Solvers and Numerical Methods for Fluid Dynamics: A Practical Introduction}.
\newblock Springer-Verlag, Berlin Heidelberg, 3rd edition.

\bibitem[{Von Neumann} and Richtmyer(1950)]{VonNeumann1950}
{Von Neumann}, J. and Richtmyer, R.D. (1950).
\newblock A method for the numerical calculation of hydrodynamic shocks.
\newblock \emph{Journal of Applied Physics}, 21(3):232--237.

\end{thebibliography}

\end{document}
